\documentclass[10pt,a4paper]{amsart}
\usepackage[margin=19mm]{geometry}
\usepackage{amsmath,amssymb,mathtools}
\usepackage[scr=rsfs,cal=euler]{mathalfa}
\usepackage{xfrac}
\usepackage[unicode,hidelinks,bookmarks=false]{hyperref}
\newcommand{\ie}{i.e.\ }
\newcommand{\cf}{cf.\ }
\newcommand{\resp}{respectively\ }
\newcommand{\Subsection}{\subsection}
\providecommand{\nobreakdash}{\nobreak\hbox{-}}
\setlength{\emergencystretch}{2em}
\multlinegap0pt
\usepackage{bm}
\usepackage{bbm}
\usepackage{dsfont}
\usepackage{xspace}
\usepackage{cancel}
\usepackage{mathtools}
\usepackage{amsthm}
\makeatletter
\@ifundefined{theorem}{\newtheorem{theorem}{Theorem}}{}
\@ifundefined{proposition}{\newtheorem{proposition}[theorem]{Proposition}}{}
\makeatother
\renewcommand{\emph}[1]{\textit{#1}}
\title[Discrete-time, discrete-state multistate Markov models from ]{Discrete-time, discrete-state multistate Markov models from the perspective of algebraic statistics}
\author{Dario Gasbarra, Kaie Kubjas, Sangita Kulathinal, Nataliia Kushnerchuk, Fatemeh Mohammadi, Etienne Sebag}
\date{Source: arXiv:2602.09619v2 (CC BY 4.0)}
\begin{document}
\maketitle

\noindent\emph{Source and licence.} Discrete-time, discrete-state multistate Markov models from the perspective of algebraic statistics. Version arXiv:2602.09619v2 (CC BY 4.0), distributed under the Creative Commons Attribution 4.0 International licence, as recorded in the arXiv licence metadata for this exact version and shown on the abstract page https://arxiv.org/abs/2602.09619v2. Full rights notice: licenses/arxiv-2602.09619-CC-BY-4.0.txt.\par\smallskip

\noindent\textit{Editorial scope.} Counted unit: the Proposition whose source label is \texttt{prop:homogeneous binomials} in the article (source0.tex lines 1286--1314, source label \texttt{prop:homogeneous binomials}), delivered together with the article's own complete original proof of it (source0.tex lines 1317--1441, one \texttt{proof} environment of 125 lines). No mathematics is written, repaired or re-derived: statement and proof are the article's own text line for line. Required context retained uncounted: none. Every definition, display and label of those lines is retained unchanged. Omitted from this package and not redistributed: 1--1285, 1315--1316, 1442--2304. Nothing inside a retained excerpt is omitted or altered: every retained body is delivered exactly as the article's lines are written, so no omission receipt of any kind accompanies this package. Citations: the retained excerpts carry no citation command at all, so no bibliography entry of the article is transcribed and no reference is redistributed. Reference adaptation: none was needed and none was made; every reference inside the delivered unit resolves inside it, so no unresolved reference or citation is produced. Printed slips kept literal: no printed word, symbol, hypothesis or number of the counted statement or of its proof was altered. Layout and class: the article's own class is \texttt{article} with packages including inputenc, babel, dsfont, enumerate, geometry, enumitem, bbding, amsmath, multirow, amsthm, amssymb, subcaption, mathrsfs, mathtools; the wrapper is the lane's pinned \texttt{amsart} editorial wrapper at 10pt on A4, which restates the theorem-like environments the retained text needs and adds the article's own macro definitions that the retained text needs (\texttt{\textbackslash emph}), with extra wrapper packages (bm, bbm, dsfont, xspace, cancel, mathtools). The delivered numbering is the unit's own. Assets: no retained span contains a figure environment, a graphics-inclusion command, a PDF-page inclusion, a table or a TikZ picture; no image file is copied into this package, the package declares no asset dependency and no image file is redistributed with it. Licence: version arXiv:2602.09619v2 of this article is distributed under the Creative Commons Attribution 4.0 International licence, as recorded in the arXiv licence metadata for this exact version and shown on the abstract page https://arxiv.org/abs/2602.09619v2. A full rights notice is delivered at licenses/arxiv-2602.09619-CC-BY-4.0.txt. Characters: every retained excerpt is pure ASCII, so the delivered engine, XeLaTeX with its default Latin Modern text font, reports no missing character.
\begin{proposition}
\label{prop:homogeneous binomials}
Let $\mathcal M$ be the $k$-th order homogeneous multistate Markov chain.
Fix two states $\psi,\eta$ and two consecutive blocks of positions in a path of length $k$
\[
\Gamma=(\gamma_1,\ldots,\gamma_k), \qquad
\Delta=(\delta_1,\ldots,\delta_k).
\]
Then the vanishing ideal $I(\mathcal M)$ contains all binomials of the form
\[
p_{I,\,\Gamma,\,\psi,\,\Delta,\,J}\,
p_{I',\,\Gamma,\,\eta,\,\Delta,\,J'}
-
p_{I,\,\Gamma,\,\eta,\,\Delta,\,J}\,
p_{I',\,\Gamma,\,\psi,\,\Delta,\,J'},
\]
where $I,I',J,J'$ are index blocks such that the concatenations
\[
I\,\Gamma\,\psi\,\Delta\,J,\quad
I'\,\Gamma\,\psi\,\Delta\,J',\quad
I\,\Gamma\,\eta\,\Delta\,J,\quad
\text{and}\quad
I'\,\Gamma\,\eta\,\Delta\,J'
\]
define paths of length $n$.
If one of the blocks $I$ or $I'$ is empty (respectively, $J$ or $J'$ is empty),
then the corresponding block $\Gamma$ (respectively, $\Delta$) 
may have length strictly smaller than $k$.
\end{proposition}

\begin{proof}
Assume $I,J,I',J'$ are nonempty (the empty-block case is identical). Fix states
$\psi,\eta$ and fixed states $\gamma_1,\dots,\gamma_k$ and $\delta_1,\dots,\delta_k$.
We first treat the interior case, i.e.\ when the symbol $\psi$ (or $\eta$) occurs at a position
$r$ with $k+1\le r,\ell\le n-k$, so that the
corresponding joint-probability entries are
\[
p_{i_1,\dots,i_{r-k-1},\gamma_1,\dots,\gamma_k,\psi,\delta_1,\dots,\delta_k,i_{r+k+1},\dots,i_n}
\quad\text{and}\quad
p_{j_1,\dots,j_{\ell-k-1},\gamma_1,\dots,\gamma_k,\psi,\delta_1,\dots,\delta_k,j_{\ell+k+1},\dots,j_n},
\]
and likewise with $\psi$ replaced by $\eta$.


\medskip

Since we are in the \emph{homogeneous} $k$-th order Markov model, the joint probability of a path
factors as
\[
p_{x_1,\dots,x_n}
=\pi_{x_1,\dots,x_k}\,\prod_{t=k+1}^{n} a_{x_{t-k}\cdots x_{t-1}x_t},
\]
where the transition parameters $a_{x_{t-k}\cdots x_{t-1}x_t}$ do \emph{not} depend on $t$.

Consider the ratio of the two path probabilities that differ only by $\psi$ versus $\eta$ at the
position $r$. By cancellation of all common factors in the above product, we get
\begin{align*}
\frac{
p_{i_1,\dots,i_{r-k-1},\gamma_1,\dots,\gamma_k,\psi,\delta_1,\dots,\delta_k,i_{r+k+1},\dots,i_n}
}{
p_{i_1,\dots,i_{r-k-1},\gamma_1,\dots,\gamma_k,\eta,\delta_1,\dots,\delta_k,i_{r+k+1},\dots,i_n}
}
&=
\frac{
a_{\gamma_1\cdots\gamma_k\psi}\,
a_{\gamma_2\cdots\gamma_k\psi\delta_1}\cdots
a_{\psi\delta_1\cdots\delta_k}
}{
a_{\gamma_1\cdots\gamma_k\eta}\,
a_{\gamma_2\cdots\gamma_k\eta\delta_1}\cdots
a_{\eta\delta_1\cdots\delta_k}
}.
\end{align*}
Exactly the same cancellation argument for the ratio at the position $\ell$ yields
\begin{align*}
\frac{
p_{j_1,\dots,j_{\ell-k-1},\gamma_1,\dots,\gamma_k,\psi,\delta_1,\dots,\delta_k,j_{\ell+k+1},\dots,j_n}
}{
p_{j_1,\dots,j_{\ell-k-1},\gamma_1,\dots,\gamma_k,\eta,\delta_1,\dots,\delta_k,j_{\ell+k+1},\dots,j_n}
}
&=
\frac{
a_{\gamma_1\cdots\gamma_k\psi}\,
a_{\gamma_2\cdots\gamma_k\psi\delta_1}\cdots
a_{\psi\delta_1\cdots\delta_k}
}{
a_{\gamma_1\cdots\gamma_k\eta}\,
a_{\gamma_2\cdots\gamma_k\eta\delta_1}\cdots
a_{\eta\delta_1\cdots\delta_k}
}.
\end{align*}
Hence, the two ratios are equal.
Equivalently, by Bayes' formula, the equality holds
if and only if
\begin{align*}
P\!\left(
X_r=\psi \mid 
X_1=i_1,\ldots,X_{r-k-1}=i_{r-k-1},\,
X_{r-k}=\gamma_1,\ldots,X_{r-1}=\gamma_k, \right. \\
\left.
X_{r+1}=\delta_1,\ldots,X_{r+k}=\delta_k,\,
X_{r+k+1}=i_{r+k+1},\ldots,X_n=i_n
\right)
\\
=
P\!\left(
X_\ell=\psi \mid 
X_1=j_1,\ldots,X_{\ell-k-1}=j_{\ell-k-1},\,
X_{\ell-k}=\gamma_1,\ldots,X_{\ell-1}=\gamma_k, \right. \\
\left.
X_{\ell+1}=\delta_1,\ldots,X_{\ell+k}=\delta_k,\,
X_{\ell+k+1}=j_{\ell+k+1},\ldots,X_n=j_n
\right)
\\
=
P\!\left(
X_r=\psi \mid
X_{r-k}=\gamma_1,\ldots,X_{r-1}=\gamma_k,\,
X_{r+1}=\delta_1,\ldots,X_{r+k}=\delta_k
\right)
\\
=
P\!\left(
X_\ell=\psi \mid
X_{\ell-k}=\gamma_1,\ldots,X_{\ell-1}=\gamma_k,\,
X_{\ell+1}=\delta_1,\ldots,X_{\ell+k}=\delta_k
\right).
\end{align*}
By the $k$-th order Markov property, conditioning on the full configuration may be
 reduced to conditioning on the $k$ nearest neighbors, 
\[
P(X_r=\psi \mid X_1,\ldots,X_{r-1},X_{r+1},\ldots,X_n)
=
P(X_r=\psi \mid X_{r-k},\ldots,X_{r-1},X_{r+1},\ldots,X_{r+k}),
\]
for $k+1\le r\le n-k$. Moreover, by homogeneity, this conditional distribution depends only on the local block
$(\gamma_1,\ldots,\gamma_k,\,\psi,\,\delta_1,\ldots,\delta_k)$ and is therefore independent
of the absolute position~$r$.
This is precisely the defining property of a homogeneous $k$-th order Markov chain.

Finally, clearing denominators yields the binomial
\[
p_{I,\,\gamma_1\ldots\gamma_k,\,\psi,\,\delta_1\ldots\delta_k,\,J}\,
p_{I',\,\gamma_1\ldots\gamma_k,\,\eta,\,\delta_1\ldots\delta_k,\,J'}
-p_{I,\,\gamma_1\ldots\gamma_k,\,\eta,\,\delta_1\ldots\delta_k,\,J}\,
p_{I',\,\gamma_1\ldots\gamma_k,\,\psi,\,\delta_1\ldots\delta_k,\,J'},
\]
which therefore belongs to $I(\mathcal{M})$.

For boundary cases (when $r\le k$ or $r>n-k$), the same cancellation argument applies:
the ratio is a product of the transition factors whose windows intersect $r$, and the set of
affected windows is again determined by the available neighbors.
Thus the ratio still depends only on the local configuration and not on the global position,
and the same cross-multiplication yields the stated binomial.
\end{proof}

\end{document}
